% id: 94-01
% book: 베이직쎈 미적분1 · 06 도함수의 활용 (3)
% section: 개념 40 방정식에의 활용
% passage: 다음 방정식의 서로 다른 실근의 개수를 구하시오.
% figure: tikz:fig-94-01
% answer: 
% answer_source: 
% variant_level: 0 (원본 전사)
% 이 파일은 build.mjs 가 items.json 에서 만든다. 고칠 때는 items.json 을 고친다.
\smash{\raisebox{1.5mm}{\dmkichul{베이직쎈 미적분1 94쪽 01번}}}
\noindent\begin{minipage}[t]{0.58\linewidth}\vspace{0pt}\raggedright $x^3-6x^2+9x-1=0$\end{minipage}\hfill\begin{minipage}[t]{0.4\linewidth}\vspace{0pt}\centering \resizebox{\ifdim\width>\linewidth\linewidth\else\width\fi}{!}{% 94-01(94쪽) — 01번 풀이 힌트 오른쪽의 그림 · $y=f(x)=x^3-6x^2+9x-1$ 의 그래프($x=1$ 에서 극대 3, $x=3$ 에서 극소 $-1$). 스캔 화질이 낮아 TikZ 로 다시 그림.
% 원문에 있는 것만 옮긴다: 곡선 · 눈금 라벨 3, $-1$($y$축) 과 1, 3($x$축) · 안내 점선 네 줄 · 라벨 O, x, y, $y=f(x)$. 원문에 없는 값($x$축 교점의 좌표)은 적지 않는다.
\begin{tikzpicture}[x=0.45cm, y=0.33cm, line width=0.6pt]
  \draw[->, >=stealth, line width=0.7pt] (-0.8,0) -- (5.3,0) node[below=2pt, font=\small] {$x$};
  \draw[->, >=stealth, line width=0.7pt] (0,-2.8) -- (0,4.7) node[left=1pt, font=\small] {$y$};
  \draw[densely dotted, line width=0.4pt] (0,3) -- (1,3);
  \draw[densely dotted, line width=0.4pt] (1,3) -- (1,0);
  \draw[densely dotted, line width=0.4pt] (0,-1) -- (3,-1);
  \draw[densely dotted, line width=0.4pt] (3,0) -- (3,-1);
  \draw plot[smooth, tension=0.7] coordinates {(-0.12,-2.17) (0,-1) (0.2,0.57) (0.4,1.7) (0.6,2.46) (0.8,2.87) (1,3) (1.2,2.89) (1.4,2.58) (1.6,2.14) (1.8,1.59) (2,1) (2.2,0.41) (2.4,-0.14) (2.6,-0.58) (2.8,-0.89) (3,-1) (3.2,-0.87) (3.4,-0.46) (3.6,0.3) (3.8,1.43) (4,2.6) (4.1,3.96)};
  \node[left=1pt, font=\small] at (0,3) {$3$};
  \node[left=1pt, font=\small] at (0,-1) {$-1$};
  \node[below=1pt, font=\small] at (1,0) {$1$};
  \node[above=1pt, font=\small] at (3,0) {$3$};
  \node[above left=-2pt and -1pt, font=\small] at (0,0) {$\pt{O}$};
  \node[font=\small] at (3.5,4.35) {$y=f(x)$};
\end{tikzpicture}
}\end{minipage}\par
\dmhint{$f(x)=x^3-6x^2+9x-1$로 놓으면\\ \hspace*{1.5em}$f'(x)=3x^2-12x+9=3(x-1)(\blank[2.4em])$\\ $f'(x)=0$에서 \qquad $x=1$ 또는 $x=\blank$\par\smallskip\begin{center}\begin{tabular}{c|c|c|c|c|c}\hline $x$ & $\cdots$ & $1$ & $\cdots$ & \blank & $\cdots$ \\ \hline $f'(x)$ & $+$ & $0$ & $-$ & $0$ & $+$ \\ \hline $f(x)$ & $\nearrow$ & $3$ & $\searrow$ & \blank & $\nearrow$ \\ \hline\end{tabular}\end{center}\smallskip\noindent 따라서 함수 $y=f(x)$의 그래프는 오른쪽 그림과 같이 $x$축과 서로 다른 세 점에서 만나므로 주어진 방정식의 서로 다른 실근의 개수는 \blank이다.}
